Fresh Graduate Teknik Informatika dengan fokus pada AI Engineering, Machine Learning, Natural Language Processing (NLP), dan Computer Vision. Memiliki pengalaman dalam pengembangan model AI menggunakan PyTorch, Hugging Face, Scikit-learn, OpenCV, serta arsitektur berbasis Transformer melalui berbagai proyek seperti multimodal CLIP + ELECTRA, sistem Retrieval-Augmented Generation (RAG) berbasis BM25, Named Entity Recognition (NER) menggunakan IndoBERT, dan klasifikasi citra dengan EfficientNet. Didukung oleh pengalaman kerja praktik di PT Pertamina Hulu Rokan.

\section*{Pendidikan}
\resumeSubheading{Institut Teknologi Sumatera— S.Kom. Teknik Informatika}{Agu 2022 - Mei 2026}{Lampung, Indonesia}{}
\begin{itemize}
  \item IPK: 3.64/4.00
  \item Fokus: AI Engineering, Machine Learning, NLP, dan Computer Vision
  \item Mata Kuliah Relevan: Machine Learning, Deep Learning, NLP, Multimodal Machine Learning
\end{itemize}

\section*{Pengalaman Kerja}
\resumeSubheading{AI Engineer Intern — PT Pertamina Hulu Rokan}{Jun 2025 - Jul 2025}{Rumbai, Pekanbaru, Indonesia}{}
\begin{itemize}
  \item Mengembangkan sistem Retrieval-Augmented Generation (RAG) dan berbasis LLM untuk tanya jawab dokumen internal menggunakan Python dan PyTorch
  \item Menerapkan teknik NLP dan membangun pipeline pengambilan dokumen untuk meningkatkan akses informasi serta mendukung inisiatif digitalisasi dokumen berbasis AI
\end{itemize}

% \section*{Pengalaman Mengajar}
\resumeSubheading{Asisten Praktikum — UPA TIK ITERA}{Agu 2023 - Sekarang}{Lampung, Indonesia}{}
\begin{itemize}
  \item \textbf{Pengenalan Komputer \& Soft1 I (PKS I):} Membimbing mahasiswa dalam pemrograman prosedural dan logika dasar (Credential ID: B/5359/IT9.E2/TA.00.00/2024)
  \item \textbf{Desain Interaksi:} Memberikan bimbingan teknis untuk pengembangan proyek UI/UX dan prinsip desain yang berpusat pada pengguna (Credential ID: 33133/IT9.3.3/TA.00.00/2025)
  \item \textbf{Sistem Operasi:} Membimbing praktikum berbasis Linux, termasuk manajemen proses, manajemen memori, dan penggunaan command line Linux
\end{itemize}

\section*{Proyek}
\resumeProject{Multimodal Deep Learning for Self-Harm Meme Detection $|$ \href{https://github.com/elsaelisa09/multimodal-meme-selfharm}{\footnotesize [GitHub]}}{Agu 2025 - Apr 2026}{Teknologi: CLIP, ELECTRA, Gradio}
\begin{itemize}
  \item Mengembangkan sistem deep learning multimodal menggabungkan CLIP dan ELECTRA untuk deteksi meme self-harm  dengan akurasi 96\%
  \item Melakukan deployment model menggunakan Gradio pada Hugging Face Spaces
\end{itemize}

\resumeProject{Multimodal City Branding Analysis $|$ \href{https://github.com/elsaelisa09/city-branding-multimodal}{\footnotesize [GitHub]}}{Mei 2026 - Sekarang}{Teknologi: Image-Text Representation Learning, Analisis Media Sosial}
\begin{itemize}
  \item Mengembangkan sistem AI multimodal untuk analisis pariwisata menggunakan klasifikasi data gambar-teks Instagram
  \item Proyek kolaborasi dengan program studi Perencanaan Wilayah dan Kota ITERA dalam riset pemodelan AI
\end{itemize}

\resumeProject{Student Face Recognition with ArcFace and SVM $|$ \href{https://github.com/elsaelisa09/student-face-recognition-arcface-svm}{\footnotesize [GitHub]}}{Jun 2026}{Teknologi: InsightFace ArcFace, Scikit-learn, OpenCV, SVM, Jupyter Notebook}
\begin{itemize}
  \item Mengembangkan sistem identifikasi wajah end-to-end untuk mengotomatisasi pengenalan 70 identitas mahasiswa dari gambar wajah yang diunggah.
  \item Membangun pipeline menggunakan augmentasi gambar offline, embedding 512 dimensi InsightFace ArcFace Buffalo\_L, dan klasifikasi SVM kernel RBF dengan evaluasi Stratified Group K-Fold yang tahan kebocoran.
  \item Memproses 1.946 sampel wajah dan mencapai akurasi 82,0\%, presisi tertimbang 85,4\%, dan F1-score tertimbang 82,4\%.
\end{itemize}

\resumeProject{Document AI ChatBot with RAG \& LLM $|$ \href{https://github.com/elsaelisa09/ChatBot-Dokumen}{\footnotesize [GitHub]}}{Jun 2025 - Jul 2025}{Teknologi: RAG, Integrasi LLM, Pencarian Semantik}
\begin{itemize}
  \item Membangun chatbot Document AI yang mengintegrasikan Retrieval-Augmented Generation dengan LLM untuk tanya jawab dokumen cerdas
  \item Mengembangkan pipeline pengambilan dan pembuatan respons kontekstual untuk pencarian semantik
\end{itemize}

\section*{Kepemimpinan \& Organisasi}
\resumeSubheading{Kepala Divisi Akademik \& Beasiswa — HMIF ITERA}{Jan 2025 – Des 2025}{Lampung, Indonesia}{}
\begin{itemize}
  \item Mengelola program akademik dan peluang beasiswa bagi mahasiswa Teknik Informatika ITERA
\end{itemize}

\resumeSubheading{Kepala Departemen Kewirausahaan — Sobat Bumi Lampung}{Okt 2025 – Sekarang}{Lampung, Indonesia}{}
\begin{itemize}
  \item Memimpin inisiatif kewirausahaan dan proyek komunitas yang berfokus pada keberlanjutan dan kesadaran lingkungan
\end{itemize}

\resumeSubheading{Sekretaris Menteri Informatika — KM ITERA}{Apr 2024 – Des 2024}{Lampung, Indonesia}{}
\begin{itemize}
  \item Mengelola dokumentasi administrasi dan koordinasi proyek IT Keluarga Mahasiswa ITERA
\end{itemize}

\resumeSubheading{Staf Divisi Publikasi \& Dokumentasi — HMIF ITERA}{Mar 2024 – Des 2024}{Lampung, Indonesia}{}
\begin{itemize}
  \item Mendokumentasikan kegiatan organisasi Himpunan Mahasiswa Informatika
\end{itemize}

\section*{Keahlian}
\begin{itemize}
  \item \textbf{AI/ML:} Machine Learning, Deep Learning, NLP, Computer Vision, LLMs, RAG, Fine-Tuning
  \item \textbf{Pengembangan AI:} Multimodal Learning, NER, Klasifikasi Gambar, Klasifikasi Audio
  \item \textbf{Framework \& Alat:} PyTorch, RunPod, Scikit-learn, Hugging Face, OpenCV, Transformers, W\&B
  \item \textbf{Pemrograman:} Python, SQL, Git, Linux
  \item \textbf{Lainnya:} Tableau, Figma
\end{itemize}

\section*{Sertifikasi}
\begin{itemize}
  \item \textbf{Program Cloud Computing AWS re/Start} - Orbit Future Academy \& Amazon Web Services (2025 - Sekarang)
  \item \textbf{Belajar Pemrograman Prosedural dengan Python} - Dicoding (2024) | Credential ID: 2VX3O4W23ZYQ
  \item \textbf{Memulai Pemrograman dengan Python} - Dicoding (2024) | Credential ID: L4PQQ8O2PO01
  \item \textbf{Belajar Dasar Structured Query Language (SQL)} - Dicoding (2024) | Credential ID: NVP77M3WVPR0
  \item \textbf{AI Praktis untuk Productivity} - Dicoding (2024) | Credential ID: 1OP82QOYLPQK
\end{itemize}

\section*{Penghargaan}
\begin{itemize}
  \item \textbf{Penerima Beasiswa Pertamina Sobat Bumi} - Pertamina Foundation (2024)
  \item \textbf{Juara 2 Kompetisi UI/UX Point Project 2.0} - HMIF ITERA (2024) | Credential ID: 313/525/SRT/HMIF/XII/2024
\end{itemize}

\section*{Bahasa \& Publikasi}
\begin{itemize}
  \item \textbf{Indonesian:} Native — strong reading, writing, and contextual evaluation
  \item \textbf{English:} Professional Working Proficiency — able to write structured AI evaluation rationales
  \item \textbf{Publikasi:} Sianturi, E. E. Y., dkk. \textit{"Transformasi Digital Layanan Posyandu."} PADIMAS, 2025 $|$ \href{https://jurnal.mdp.ac.id/index.php/padimas/article/view/12206}{\footnotesize [Paper Lengkap]}
\end{itemize}
